\documentclass[a4paper,11pt]{article}
\def\email#1{{\tt#1}}

\def\N{\mbox{\rm{I}\hspace{-.15em}\rm{N}}}
\def\Z{\mbox{\rm{Z}\hspace{-.28em}\rm{Z}}}
\def\RR{\mbox{\rm{I}\hspace{-.15em}\rm{R}}}
\def\C{\mbox{\rm{l}\hspace{-.50em}\rm{C}}}
\def\Q{\mbox{\rm{l}\hspace{-.50em}\rm{Q}}}
\def\F{\mbox{\rm{I}\hspace{-.20em}\rm{F}}}

\usepackage{amssymb,amsmath}
\usepackage{url}
\usepackage{fullpage}
\begin{document}

\newtheorem{theo}{Theorem}
\newtheorem{coro}{Corollary}[theo]
\newtheorem{lemme}{Lemma}
\renewcommand{\thecoro}{\thetheo.\arabic{coro}}



\parindent=0cm


\section*{\center IN100 Initiation au Python\\
Contrôle continu 2}

\vspace{-0.3 cm}

\begin{table}[h!]
\centering
\renewcommand{\arraystretch}{1.4} % augmente la hauteur des lignes
\begin{tabular}{|p{10cm}|p{3cm}|p{2cm}|}
\hline
\textbf{Étudiant :} & \textbf{Heure Début }  & \textbf{Note (/5) }  \\
& & \\
\hline
\textbf{Correcteur :} & \textbf{Heure Fin } & \\
& & \\
\hline

\hline
\end{tabular}
\end{table}
\vspace{-0.4 cm}
\paragraph*{Consignes:} Vous avez 20 minutes pour coder la solution. Seules les types et les instructions vues en cours sont autorisées. Tout usage d'une fonction 
venant d'une bibliothèque qui n'a pas été vue en cours ne sera pris en compte. 

\vspace{-0.2 cm}
\vspace{-0.2 cm}
\paragraph*{Sujet:}

On considère une matrice carrée.  
On dit que cette matrice est \textbf{diagonale dominante} si, pour chaque ligne $i$, on a :

\[
|a_{ii}| \ge \sum_{j \neq i} |a_{ij}|.
\]

Cette condition signifie que, sur chaque ligne, l'élément situé sur la diagonale possède une valeur absolue supérieure ou égale à la somme des valeurs absolues des autres éléments de la même ligne. 



\textbf{Exemple :}  
Considérons la matrice suivante :
\[
\begin{bmatrix}
5 & 1 & 2 \\
0 & 7 & 1 \\
2 & 0 & 3
\end{bmatrix}
\]

Pour la ligne 0 :
\[
|5| \ge |1| + |2|
\]

Pour la ligne 1 :
\[
|7| \ge |0| + |1|
\]

Pour la ligne 2 :
\[
|3| \ge |2| + |0|
\]

Cette matrice est donc diagonale dominante.

\vspace{-0.3 cm}


\begin{enumerate}
  \item[1.] Écrire une fonction qui détermine si une ligne donnée d’une matrice vérifie la condition de dominance diagonale.
  
  \item[2.] Écrire une fonction qui indique si une matrice entière est diagonale dominante.
\end{enumerate}


\vspace{-0.4 cm}
\begin{table}[h!]
\centering
\renewcommand{\arraystretch}{1.4} % augmente la hauteur des lignes
\begin{tabular}{|p{15cm}|}
\hline
\textbf{Remarques étudiant :}    \\
 \\
\hline
\textbf{Remarques correcteur :}  \\
 \\
\hline
\end{tabular}
\end{table}

\newpage


\section*{\center IN100 Initiation au Python\\
Contrôle continu 2}

\vspace{-0.3 cm}

\begin{table}[h!]
\centering
\renewcommand{\arraystretch}{1.4} % augmente la hauteur des lignes
\begin{tabular}{|p{10cm}|p{3cm}|p{2cm}|}
\hline
\textbf{Étudiant :} & \textbf{Heure Début }  & \textbf{Note (/5) }  \\
& & \\
\hline
\textbf{Correcteur :} & \textbf{Heure Fin } & \\
& & \\
\hline

\hline
\end{tabular}
\end{table}
\vspace{-0.4 cm}
\paragraph*{Consignes:} Vous avez 20 minutes pour coder la solution. Seules les types et les instructions vues en cours sont autorisées. Tout usage d'une fonction 
venant d'une bibliothèque qui n'a pas été vue en cours ne sera pris en compte. 



\paragraph*{Sujet:}



Créez une liste imbriquée (double liste)  de taille 7, où chaque sous-liste est aussi de taille 7 et où les éléments sont initialisés à 0. 
\begin{enumerate}
  \item[1.] Donner une fonction qui place un 1 aléatoirement dans une  liste imbriquée.   
  \item[2.] Une liste imbriquée peut être  vue comme une grille carrée, ou une matrice. 
\'Ecrire  une fonction qui déplace le 1 dans la liste aléatoirement d'une position vers le haut ou vers le bas ou vers la gauche ou vers la droite. % Enfin, faites une boucle qui applique cette dernière fonction jusqu'à ce que le 1 touche un bord de la grille. 
\end{enumerate}

\textbf{Exemples}:
\noindent
\begin{itemize}
\item Si $a_{23}=1$  alors le déplacement vers le bas est  $a_{33}=1$.  \\
\item Si $a_{34}=1$ alors le déplacement vers la gauche est $a_{33}=1$ \\
\end{itemize} 

\`A chaque déplacement, il faudra faire attention de ne pas déborder. 


\vspace{0.2 cm}
Pour générer aléatoirement des nombres entiers, nous rappelons qu'on peut utiliser la fonction \texttt{random.randint(a,b)} qui renvoie un nombre entier entre $a$ et $b$. 

\vspace{0.2 cm}
Exemple d'usage:
\vspace{-0.2 cm}
\begin{verbatim}
	import random
	
	print(random.randint(2,15))
\end{verbatim}




\vspace{0.5 cm}
\begin{table}[h!]
\centering
\renewcommand{\arraystretch}{1.4} % augmente la hauteur des lignes
\begin{tabular}{|p{15cm}|}
\hline
\textbf{Remarques étudiant :}    \\
 \\
\hline
\textbf{Remarques correcteur :}  \\
 \\
\hline
\end{tabular}
\end{table}

\newpage

\section*{\center IN100 Initiation au Python\\
Contrôle continu 2}

\vspace{-0.3 cm}

\begin{table}[h!]
\centering
\renewcommand{\arraystretch}{1.4} % augmente la hauteur des lignes
\begin{tabular}{|p{10cm}|p{3cm}|p{2cm}|}
\hline
\textbf{Étudiant :} & \textbf{Heure Début }  & \textbf{Note (/5) }  \\
& & \\
\hline
\textbf{Correcteur :} & \textbf{Heure Fin } & \\
& & \\
\hline

\hline
\end{tabular}
\end{table}
\vspace{-0.4 cm}
\paragraph*{Consignes:} Vous avez 20 minutes pour coder la solution. Seules les types et les instructions vues en cours sont autorisées. Tout usage d'une fonction 
venant d'une bibliothèque qui n'a pas été vue en cours ne sera pris en compte. 


\vspace{0.2 cm}
\paragraph*{Sujet:}


\begin{enumerate}
		
		\item Ecrire la fonction \texttt{moyenne} qui prend pour argument une liste d'entiers et qui retourne la moyenne de ces nombres.
		\item Initialiser un dictionnaire qui contient les liste des notes obtenues pendant l'année par les étudiants d'une licence: 
		\begin{verbatim}
		base =  {"Leonard": [12, 8, 11]],  "Sophie": [14, 12, 16, 8, 13]}
		\end{verbatim}
		
		\noindent \'Ecrire la fonction \texttt{calcul\_moyenne} qui prend pour argument un dictionnaire \texttt{base}   et le nom d'un étudiant. Cette fonction doit retourner la moyenne de l'étudiant si celui-ci est présent dans la base, et -1 sinon.
		Tester le programme.
		
	\end{enumerate}

\vspace{1 cm}
\begin{table}[h!]
\centering
\renewcommand{\arraystretch}{1.4} % augmente la hauteur des lignes
\begin{tabular}{|p{15cm}|}
\hline
\textbf{Remarques étudiant :}    \\
 \\
\hline
\textbf{Remarques correcteur :}  \\
 \\
\hline
\end{tabular}
\end{table}




\newpage

\section*{\center IN100 Initiation au Python\\
Contrôle continu 2}

\vspace{0.3 cm}

\begin{table}[h!]
\centering
\renewcommand{\arraystretch}{1.4} % augmente la hauteur des lignes
\begin{tabular}{|p{10cm}|p{3cm}|p{2cm}|}
\hline
\textbf{Étudiant :} & \textbf{Heure Début }  & \textbf{Note (/5) }  \\
& & \\
\hline
\textbf{Correcteur :} & \textbf{Heure Fin } & \\
& & \\
\hline

\hline
\end{tabular}
\end{table}
\vspace{0.4 cm}
\paragraph*{Consignes:} Vous avez 20 minutes pour coder la solution. Seules les types et les instructions vues en cours sont autorisées. Tout usage d'une fonction 
venant d'une bibliothèque qui n'a pas été vue en cours ne sera pris en compte. 

\vspace{-0.2 cm}
\vspace{-0.2 cm}
\paragraph*{Sujet:}

On considère une suite numérique $(u_n)$ définie par une valeur initiale $u_0 = a$, et par la relation de récurrence suivante, valable pour tout entier $n \ge 0$ :
\[
u_{n+1} = u_n + 3.
\]

Cette relation signifie que chaque terme de la suite est obtenu en ajoutant 3 au terme précédent.

\vspace{0.5 cm}

\textbf{Exemple :}  
Si $u_0 = 2$, alors :
\[
u_1 = 2 + 3 = 5
\]
\[
u_2 = 5 + 3 = 8
\]
\[
u_3 = 8 + 3 = 11
\]

La suite commence donc par :
\[
(2,\; 5,\; 8,\; 11,\; \ldots)
\]



\begin{enumerate}
\item[1.] Écrire une fonction \texttt{u(u0)} qui prend en entrée $u_0$ et qui renvoie  le terme $u_n$ de la suite.
\item[2.] Écrire une fonction \texttt{moyenne(u0,n)} qui prend en entrée $u_0$ et $n$ et qui renvoie   qui renvoie la moyenne des $n$ premières termes de la suite. 
\end{enumerate}


\vspace{1 cm}
\begin{table}[h!]
\centering
\renewcommand{\arraystretch}{1.4} % augmente la hauteur des lignes
\begin{tabular}{|p{15cm}|}
\hline
\textbf{Remarques étudiant:}    \\
 \\
\hline
\textbf{Remarques correcteur:}  \\
 \\
\hline
\end{tabular}
\end{table}

\newpage



\newpage

\section*{\center IN100 Initiation au Python\\
Contrôle continu 2}

\vspace{0.3 cm}

\begin{table}[h!]
\centering
\renewcommand{\arraystretch}{1.4} % augmente la hauteur des lignes
\begin{tabular}{|p{10cm}|p{3cm}|p{2cm}|}
\hline
\textbf{Étudiant :} & \textbf{Heure Début }  & \textbf{Note (/5) }  \\
& & \\
\hline
\textbf{Correcteur :} & \textbf{Heure Fin } & \\
& & \\
\hline

\hline
\end{tabular}
\end{table}
\vspace{0.4 cm}
\paragraph*{Consignes:} Vous avez 20 minutes pour coder la solution. Seules les types et les instructions vues en cours sont autorisées. Tout usage d'une fonction 
venant d'une bibliothèque qui n'a pas été vue en cours ne sera pris en compte. 


\paragraph*{Sujet:}


On considère une base de données simplifiée d’étudiants, représentée par un dictionnaire Python.  
Chaque clé du dictionnaire est le nom de famille d’un étudiant, et la valeur associée est une chaîne de caractère qui représente le  numéro de téléphone de l’étudiant.

Par exemple, la base de données suivante est donnée :

\begin{center}
\texttt{base = \{"Euler" :  "01.43.64.00.00", "Germain" : "Sophie", "01.43.64.11.11" \}}.
\end{center}


\begin{enumerate}
   
    \item Écrire la fonction Python \texttt{quel\_numero(base, nom)} qui prend en argument une base de données
    \texttt{base} et un nom de famille \texttt{nom}.  
    Cette fonction doit retourner le numéro de téléphone de l’étudiant si le nom est présent dans la base,
    et \texttt{-1} sinon.

    \item Écrire la fonction Python \texttt{rajout\_numero(base, numero)} qui prend en argument un dictionnaire 
    \texttt{base}, un nom \texttt{nom} et un numéro de téléphone \texttt{numero}.  
    Cette fonction vérifie que la chaine de caracteres numero a le bon format (10 chiffres, XX.XX.XX.XX.XX) 
    et rajoute la paire \texttt{nom, numero} dans le dictionnaire. 
    
    
    \item Tester les deux fonctions en les appelant successivement sur la base de données dans l'exemple. 
    \end{enumerate}

\vspace{1 cm}
\begin{table}[h!]
\centering
\renewcommand{\arraystretch}{1.4} % augmente la hauteur des lignes
\begin{tabular}{|p{15cm}|}
\hline
\textbf{Remarques étudiant:}    \\
 \\
\hline
\textbf{Remarques correcteur:}  \\
 \\
\hline
\end{tabular}
\end{table}



\end{document}
